\section{相关工作}
\label{sec:related_work}

\subsection{VLA 模型与泛化挑战}

近年来，机器人学习发生了范式转变：从专门的策略网络走向通用 VLA 模型。
$\pi_0$~\cite{pi0} 和 OpenVLA~\cite{openvla} 等先驱工作利用海量网络规模数据集（例如 Open-X Embodiment~\cite{openx}）实现端到端机器人控制，将视觉编码器（例如 SigLIP~\cite{siglip}）与大语言模型相结合，把开放词汇指令转换为精确的连续动作。
然而，通过微调（例如 X-VLA~\cite{xvla}）使这些模型适应 OOD 场景会产生高昂的计算成本，因此需要无需更新权重、参数高效且可即插即用的泛化方法。

\subsection{机器人领域的 RAG 与上下文学习}

为在不微调的情况下提高适应能力，研究者已将 RAG 和上下文学习（ICL）引入机器人领域。
RICL~\cite{ricl} 和 MemER~\cite{memer} 等框架通过向量数据库检索相关历史示范，并把它们作为上下文提示附加到策略执行中，从而通过最近邻模仿适应新任务。
这些方法虽然能有效利用外部记忆，却只关注成功轨迹，忽略了失败经历的诊断价值。这说明我们需要一种同时利用成功与失败进行对比式指导的双记忆结构。

\subsection{上下文内自适应智能体框架}

除被动检索轨迹外，近期研究还探索了让 LLM 和 VLM 作为具有外部工具访问能力的高层规划器。
Code-as-Policies~\cite{cap} 为层级任务生成可执行 Python 代码，VoxPoser~\cite{voxposer} 将指令落地为用于运动规划的三维空间价值图，而 MOKA~\cite{moka} 则通过视觉提示在图像上标注可供性。
这些框架虽然能把复杂指令分解为可在上下文中自适应的可执行原语，但缺少闭环适应所需的经验记忆与动态纠错能力。

\endinput
